%%%%%%%%%%%%%%%%%%%%%%%%%%%%%%%%%%%%%%%%%%%%%%%%%%%%%%%%%%%%%%%%
\section{实验设计}
%%%%%%%%%%%%%%%%%%%%%%%%%%%%%%%%%%%%%%%%%%%%%%%%%%%%%%%%%%%%%%%%

\subsection{实验对象}

本实验选取五款国内外主流大语言模型作为AI合成被试，以覆盖差异化的技术架构、推理能力和安全对齐策略。
模型选择遵循三个标准：其一，在中文语境中具有广泛的用户基础或学术关注度；其二，技术架构和对齐策略具有可辨识的差异，以确保模型间异质性假设的可检验性；其三，提供稳定的API访问接口以支持大规模重复实验。
五款模型的基本信息如\cref{tab:models}所示。

\begin{table}[tb]
  \caption{实验所用大语言模型的基本信息}
  \label{tab:models}
  \centering
  \footnotesize
  \setlength{\tabcolsep}{4pt}
  \begin{tabular}{p{0.68in}p{0.52in}p{1.8in}}
    \toprule
    模型名称 & 开发方 & 纳入理由 \\
    \midrule
    DeepSeek-V4-Flash & 深度求索 & DeepSeek系列轻量版本，用于检验同厂商不同推理配置的差异 \\
    DeepSeek-V4-Pro & 深度求索 & DeepSeek系列高阶版本，用于与Flash版本形成对照 \\
    Grok-4 & xAI & 国际主流闭源模型之一，用于扩展跨厂商比较 \\
    Kimi-k2 & 月之暗面 & 中文语境中常用的国内模型，用于检验中文任务下的模型差异 \\
    GPT-4.1 & OpenAI & 国际主流闭源模型之一，用于提供跨生态系统参照 \\
    \bottomrule
  \end{tabular}
\end{table}

每款模型在每个实验条件下独立运行50次重复实验，每次实验为独立API会话，无上下文共享。
重复实验采用各模型API的默认采样参数；本研究未系统操纵Temperature，因此不将其作为因果解释变量。
全部实验条件共6组，主实验总数据量为5（模型）$\times$ 6（条件）$\times$ 50（重复）= 1500组有效观测。
此外，本研究对DeepSeek-V4-Flash和DeepSeek-V4-Pro额外采集600组深度思考模式对照样本，并设置逐题独立调用的序列效应对照组以评估题目呈现方式的影响。

\begin{table}[tb]
  \caption{API调用与实验记录口径}
  \label{tab:api_protocol}
  \centering
  \scriptsize
  \setlength{\tabcolsep}{3.5pt}
  \begin{tabular}{p{0.68in}p{0.95in}p{1.22in}}
    \toprule
    模型 & 主要实验记录 & 参数与复现说明 \\
    \midrule
    DeepSeek-V4-Flash & 主实验50次/条件；另有thinking开启对照 & 默认采样参数；thinking开关作为附加实验因素记录 \\
    DeepSeek-V4-Pro & 主实验50次/条件；另有thinking开启对照 & 默认采样参数；thinking开关作为附加实验因素记录 \\
    Grok-4 & 主实验50次/条件；逐题独立对照 & 默认采样参数；每次调用为无状态会话 \\
    Kimi-k2 & 主实验50次/条件；逐题独立对照 & 默认采样参数；每次调用为无状态会话 \\
    GPT-4.1 & 主实验50次/条件；逐题独立对照 & 默认采样参数；每次调用为无状态会话 \\
    \bottomrule
  \end{tabular}
  \vspace{0.25em}

  \parbox{0.96\linewidth}{\footnotesize 注：原始数据保存了模型、条件、thinking状态、选择序列、理由文本和解析状态；各平台默认Temperature的精确数值未在原始返回中保存，因此本文仅将其作为研究局限讨论。}
\end{table}

\subsection{实验场景与任务}

实验场景设定为金融投资决策情境。
这一场景选择基于两重考量：其一，金融代理是LLM当前最广泛、最高风险的实际应用场景之一；其二，金额刺激便于构造精确的期望值计算和风险中性基准，满足实验设计的内部效度要求。
与开放式投资建议相比，成对彩票选择牺牲了一部分真实场景复杂性，但换来了更清晰的因果识别。
在本研究中，我们并不要求模型预测市场走势，而是要求它在已知概率和金额下选择安全或风险选项。
这样可以将“信息预测能力”与“风险承担倾向”区分开来，避免模型因为对某只股票、某个行业或某段市场叙事的知识差异而影响实验结果。

实验任务采用Holt与Laury的经典风险偏好测量范式~\cite{Holt2002Risk}。
该范式通过一系列成对彩票选择（Paired Lottery Choice）测量决策者的风险偏好，因其结构简洁、理论基准明确而被广泛应用于实验经济学研究。
具体而言，被试需要在10道选择题中，每道题在安全选项（选项A）和风险选项（选项B）之间做出选择。
从第1题至第10题，安全选项的高收益概率从10\%逐步递增至100\%，使得风险中性的决策者应在第4题至第5题之间从选项A切换至选项B。
转换点越靠后，表明决策者越偏好安全选项，风险规避程度越高。

\Cref{tab:hl_gain}展示了收益域中使用的成对彩票选择量表。
在损失域中，所有结果通过线性变换转换为损失金额：损失域金额 = 385 $-$ 收益域金额。
该变换保留了选项A的波动范围（40元）和选项B的波动范围（375元）不变，确保两个域在每一题上的期望值优势差距完全相同，风险中性转折点均落在第4题与第5题之间。
因此，收益域和损失域之间的差异主要来自结果框架本身，而不是来自期望值、波动幅度或转折点位置的重新设计。
这一点对于检验前景理论中的反射效应尤其重要：如果损失域中模型更保守或更冒险，我们可以更有把握地将其解释为损失语义和代理责任线索的作用，而不是量表数值结构改变造成的机械差异。

\begin{table*}[tb]
  \caption{收益域Holt-Laury成对彩票选择量表}
  \label{tab:hl_gain}
  \centering
  \footnotesize
  \setlength{\tabcolsep}{8pt}
  \begin{tabular}{cccccc}
    \toprule
    题号 & 选项A (安全) & 期望值 A & 选项B (风险) & 期望值 B & 期望值差 (A$-$B) \\
    \midrule
    1 & 10\%获200元，90\%获160元 & 164.0 & 10\%获385元，90\%获10元 & 37.7 & +126.3 \\
    2 & 20\%获200元，80\%获160元 & 168.0 & 20\%获385元，80\%获10元 & 84.4 & +83.6 \\
    3 & 30\%获200元，70\%获160元 & 172.0 & 30\%获385元，70\%获10元 & 131.1 & +40.9 \\
    4 & 40\%获200元，60\%获160元 & 176.0 & 40\%获385元，60\%获10元 & 177.8 & $-$1.8 \\
    5 & 50\%获200元，50\%获160元 & 180.0 & 50\%获385元，50\%获10元 & 197.5 & $-$17.5 \\
    6 & 60\%获200元，40\%获160元 & 184.0 & 60\%获385元，40\%获10元 & 235.0 & $-$51.0 \\
    7 & 70\%获200元，30\%获160元 & 188.0 & 70\%获385元，30\%获10元 & 272.5 & $-$84.5 \\
    8 & 80\%获200元，20\%获160元 & 192.0 & 80\%获385元，20\%获10元 & 310.0 & $-$118.0 \\
    9 & 90\%获200元，10\%获160元 & 196.0 & 90\%获385元，10\%获10元 & 347.5 & $-$151.5 \\
    10 & 100\%获200元，0\%获160元 & 200.0 & 100\%获385元，0\%获10元 & 385.0 & $-$185.0 \\
    \bottomrule
  \end{tabular}
\end{table*}

\subsection{实验条件与实验设计}

实验采用2（损益域：收益域 vs.\ 损失域）$\times$ 3（决策角色与问责性：为自己决策 vs.\ 替他人决策且不可观察 vs.\ 替他人决策且可观察）的被试间因子设计，共构成6组实验条件。
\Cref{tab:design}展示了完整的实验设计矩阵和各条件对应的假设检验目标。
采用被试间设计的原因在于，LLM对话具有明显的上下文记忆和题目顺序敏感性。
如果同一会话中连续经历多个角色条件，模型可能将前一条件的责任线索带入后一条件，从而污染角色操纵。
因此，本文将每一次模型-条件调用视为独立会话，使Self、Agent-Invisible和Agent-Visible之间的差异尽可能来自实验操纵本身。

\begin{table*}[tb]
  \caption{2 $\times$ 3 实验设计矩阵与假设检验目标}
  \label{tab:design}
  \centering
  \footnotesize
  \setlength{\tabcolsep}{6pt}
  \begin{tabular}{cccccl}
    \toprule
    条件 & 损益域 & 决策角色 & 问责性 & Prompt角色设定 & 核心检验目标 \\
    \midrule
    1 & 收益 & 为自己 & 无 & 以第一人称为自己做投资决策 & 基线：LLM自身风险偏好 \\
    2 & 收益 & 替他人 & 不可观察 & 金融投资经理，委托人不知具体选择 & H1A/H1B：纯粹代理效应 \\
    3 & 收益 & 替他人 & 可观察 & 金融投资经理，委托人完全观察决策 & H2：问责调节效应（收益域）\\
    4 & 损失 & 为自己 & 无 & 以第一人称为自己做损失管理决策 & 基线：LLM损失域风险偏好 \\
    5 & 损失 & 替他人 & 不可观察 & 金融投资经理，委托人不知具体选择 & H1A/H1B：损失域代理效应 \\
    6 & 损失 & 替他人 & 可观察 & 金融投资经理，委托人完全观察决策 & H2：问责调节效应（损失域）\\
    \bottomrule
  \end{tabular}
\end{table*}

六组条件的提示词设计遵循统一的结构模板：系统角色设定段、任务说明段和输出格式约束段。
角色设定段在不同条件间变化以操纵决策角色和问责性感知，任务说明段和输出格式约束段保持恒定以确保刺激材料的内部一致性。
具体而言，条件的系统角色设定以第一人称设定决策收益或损失的归属对象；条件2和条件3设定LLM为专业的金融投资经理，并分别告知委托人不可观察或可观察其决策过程；条件4至条件6在损失域中平行设置对应的角色与问责性组合。
各条件中的Holt-Laury彩票题目内容完全相同，仅在金额标签的表述上因损益域而存在收益和损失的区别。
这种设计使三个关键变量分别具有清晰的操作化含义：损益域由金额叙述中的收益/损失框架体现，代理角色由收益或损失归属对象体现，问责性由委托人能否观察决策过程体现。
除这些操纵信息外，各条件的题目数量、概率、金额、输出格式和解析规则均保持一致，以降低无关文本差异对结果的影响。

\subsection{实验流程}

实验按照以下六个步骤严格执行，以确保数据的可复现性和实验流程的标准化。

\textbf{步骤一：提示词标准化构建。}
针对6组实验条件分别编写标准化的系统Prompt和用户Prompt。
系统Prompt明确LLM的决策角色、损益归属和问责性信息；用户Prompt包含10道Holt-Laury彩票选择题，并附带输出格式要求。
所有Prompt经过预实验验证，确保LLM能够正确理解任务要求并生成可解析的回答格式。

\textbf{步骤二：API调用与数据采集。}
通过各模型的官方API接口，向每个模型在每个条件下发送独立的会话请求。
每次API调用为一次独立的无状态会话，不携带前次实验的任何上下文信息。
要求LLM在每道题的选择前输出其决策理由与思考过程（思维链），最后以标准化的标签格式给出最终选择。

\textbf{步骤三：重复实验。}
每款模型在每个条件下完成50次独立重复实验。
Temperature参数采用各模型API默认值，本文不对其进行系统操纵。
重复实验的目的在于通过多次采样捕捉LLM在给定Prompt下的回答分布特征，克服单次回答的随机波动。

\textbf{步骤四：序列效应控制。}
为评估一次会话中连续呈现10道题的序列位置是否对LLM的选择产生系统性影响（例如前序题目中的选项偏好通过上下文传播影响后续选择），设置序列效应对照组。
在对照组中，10道题以10次独立的API调用分别呈现，每次仅包含一道题，消除题目间的上下文传导路径。
通过对比连续呈现与独立呈现两种方式下的选择分布，检验序列效应是否存在及其影响量级。

\textbf{步骤五：跨模型重复。}
依次对DeepSeek-V4-Flash、DeepSeek-V4-Pro、Grok-4、Kimi-k2和GPT-4.1五款模型执行步骤一至步骤四的全部实验流程，确保实验程序的跨模型一致性。

\textbf{步骤六：数据清洗与编码。}
从API返回的文本中提取结构化的选择数据和思维链文本。
对于未能按标准格式输出的回答（如使用了非标准标签），由两名编码员独立判断其选择意图，编码员间信度需达到Cohen's $\kappa \geq 0.90$方可纳入正式分析。
实际主实验中，所有返回均可自动解析为预设JSON字段，因此未发生因格式错误导致的样本剔除。

\subsection{变量定义与测量}

\textbf{因变量一：安全选择数量（$N_{\text{safe}}$）。}
定义为每轮实验中决策者在10道题中选择安全选项（选项A）的总题数，取值范围为0至10。
$N_{\text{safe}}$越高，表明决策者的风险规避程度越强。
以风险中性决策者为基准，其在收益域和损失域中的$N_{\text{safe}}$均应等于4。

\textbf{因变量二：转换点（Switching Point，$S$）。}
定义为决策者从选项A切换至选项B的题号位置，取值范围为1至11（$S = 11$表示在所有10道题中始终选择选项A，即极度风险规避）。
对于存在多次切换（不一致选择）的回答模式，采用两种补充指标：其一，以最后一次切换位置作为有效转换点；其二，记录总切换次数作为回答一致性的辅助度量。

\textbf{因变量三：思维链风险归因。}
对LLM在每道题选择前输出的思维链文本进行结构化编码，统计其决策理由中提及的关键主题类别，包括期望值计算、对委托人责任的提及、对观察者的考量、对后悔的预期以及对安全性的强调。
每个类别的提及频次作为分析LLM决策机制的行为痕迹数据。

\textbf{自变量与控制变量。}
自变量包括损益域（收益/损失）、决策角色（为自己/替他人）和问责性（不可观察/可观察）。
控制变量包括模型身份（五款模型）和实验会话编号；Temperature采用平台默认值，未进行系统操纵。

\subsection{数据分析方法}

数据分析按照从整体到局部、从定量到定性的递进逻辑展开，分为以下五个层次。

\textbf{第一层：描述性统计与可视化。}
计算各实验条件下每款模型的安全选择数量和转换点的均值、标准差、中位数和四分位距。
绘制分组箱线图和核密度估计曲线，直观呈现不同条件下LLM风险偏好的分布特征与离群模式。

\textbf{第二层：主效应检验（H1A与H1B的检验）。}
以安全选择数量$N_{\text{safe}}$为因变量，以决策角色（为自己vs.\ 替他人）为自变量进行组间比较。
由于H1A与H1B主要针对收益域中的代理角色偏移提出，本文将收益域内Self、Agent-Invisible和Agent-Visible的比较作为主检验，采用以均值差为统计量的双侧置换检验分别评估Self与两类Agent条件之间的差异。
随后将收益域和损失域合并，采用Mann-Whitney U检验比较Self条件与合并后的Agent条件，作为代理角色总体效应的补充检验。
$N_{\text{safe}}$为0至10之间的有界计数变量，且不同模型的分布形态明显偏离正态，因此非参数检验比普通均值$t$检验更合适。
若替他人条件下$N_{\text{safe}}$显著高于为自己条件，则H1A（算法谨慎偏移假设）获得支持。
若替他人条件下$N_{\text{safe}}$显著低于为自己条件，则H1B（道德风险偏移假设）获得支持。

\textbf{第三层：问责调节效应检验（H2的检验）。}
以安全选择数量$N_{\text{safe}}$为因变量，首先在各模型内部比较Self、Agent-Invisible和Agent-Visible三类角色条件，使用Kruskal-Wallis H检验评估问责性相关条件之间是否存在总体差异。
为更直接检验H2所要求的损益域$\times$问责性模式，本文进一步在Agent条件内部构造差中差指标：
\[
  \begin{aligned}
  \text{DiD}
  =(\bar N_{L,V}-\bar N_{L,I})
  -(\bar N_{G,V}-\bar N_{G,I}).
  \end{aligned}
\]
其中$L/G$分别表示损失域/收益域，$V/I$分别表示可观察/不可观察；正的DiD表示可观察性带来的谨慎化在损失域中更强。
本文采用固定随机种子的双侧置换检验评估该差中差统计量，以避免仅凭分组图形推断交互作用。

\textbf{第四层：模型间异质性检验（H3的检验）。}
以安全选择数量$N_{\text{safe}}$为因变量，以模型身份为分组变量，采用Kruskal-Wallis H检验检验五款模型之间的总体差异。
在显著性检验之外，本文结合均值、标准差、分布图和风险分类图，对不同模型在整卷作答范式下的行为模式进行归纳，识别整卷稳定型、极度保守型和情境敏感型等行为原型。

\textbf{第五层：思维链文本分析。}
对LLM输出的思维链文本进行量化内容分析。
本文基于预定义关键词词典统计三类认知焦点的出现比例：期望值与计算、客户与问责、风险与损失。
随后比较不同角色条件、损益域和模型之间的关键词提及率，并通过可视化展示模型在代理情境下的注意力焦点转移。
该分析旨在为行为层面的风险偏移提供认知层面的机制解释，回答LLM在代理角色下究竟是因为考虑了哪些因素而改变了风险偏好。

上述五层分析构成了从行为指标到认知机制、从整体效应到边界条件的完整分析路径。
所有统计分析在Python 3.x环境下完成；描述性统计、非参数检验、置换检验和可视化脚本均随附录提交。
显著性水平统一设定为$\alpha = 0.05$，多重比较中采用Benjamini-Hochberg假发现率校正。
